\renewcommand{\Anno}{2002}

\renewcommand{\NumeroProblema}{2}

\renewcommand{\Testo}{Lungo un corridoio della Scuola Sant'Anna vi
sono $n$ aule.  Ognuna di esse \`{e} contrassegnata con un numero
intero compreso fra 1 ed $n$ (naturalmente aule differenti sono
contrassegnate da numeri differenti).  Dario osserva che la somma dei
numeri delle aule sul lato destro coincide con la somma dei numeri
delle aule sul lato sinistro.
	
Tornando l'anno successivo Dario nota che, in seguito a lavori di
ristrutturazione, un'aula \`{e} stata spostata da uno dei due lati
all'altro, mentre tutte le altre hanno conservato la loro
collocazione.  Inoltre tutte le aule, compresa quella spostata, hanno
conservato la loro numerazione.  Dario nota poi che ora la somma dei
numeri delle aule su uno dei due lati \`{e} il doppio della somma dei
numeri delle aule sull'altro lato.
	
Determinare per quali valori di $n$ si pu\`{o} avere questa
situazione.}

\renewcommand{\Risposta}{I possibili valori di $n$ sono 3, 8, 11.}

\renewcommand{\Aiutino}{Quanto vale all'inizio la somma dei numeri
delle aule su ciascuno dei 2 corridoi?  Sia ora $k$ l'aula che si
sposta: calcolando le nuove somme, ed imponendo che una sia il doppio
dell'altra, si ottiene che $k=n(n+1)/12$.  Ricordando che $k$ deve
essere intero e minore od uguale a $n$, si trovano i possibili valori
di $k$.  Per tali valori occorre poi esibire una disposizione delle
aule che soddisfi l'ipotesi.}

\renewcommand{\SoluzioneUfficiale}{\TitoloSoluzione La somma dei 
numeri di tutte le aule vale
$$1+2+\ldots+n=\frac{n(n+1)}{2},$$
dunque all'inizio la somma dei numeri delle aule su ciascuno dei 2 
lati del corridoio vale
$$\frac{n(n+1)}{4}.$$

Sia ora $k$ il numero dell'aula che si sposta. Al termine dei lavori 
le somme dei numeri delle aule sui 2 lati del corridoio saranno
$$\frac{n(n+1)}{4}-k
\hspace{3em}
\mbox{e}
\hspace{3em}
\frac{n(n+1)}{4}+k.$$

Imponendo ora che la seconda somma sia il doppio della prima, e 
semplificando, otteniamo che
\begin{equation}
	k=\frac{n(n+1)}{12}
	\label{eq:aulefranzoni}
\end{equation}

Ora facciamo due considerazioni.
\begin{itemize}
	\item  Dovendo essere $k\leq n$, sostituendo la 
	(\ref{eq:aulefranzoni}) otteniamo la disequazione $n(n-11)\leq 0$, da 
	cui $n\leq 11$.

	\item  Dovendo $k$ essere intero, dalla (\ref{eq:aulefranzoni}) 
	ricaviamo che $n(n+1)$ deve essere divisibile per 12. Per verifica 
	diretta, gli unici valori interi di $n$, compresi tra 1 e 11, per 
	cui ci\`{o} accade sono $n=3$ (da cui $k=1$), $n=8$ (da cui $k=6$) 
	e $n=11$ (da cui $k=11$).
\end{itemize}

Non ci resta ora che fornire, per tali valori di $n$, degli esempi di 
disposizione e spostamento delle aule che verificano le condizioni 
richieste (i corrispondenti valori di $k$ ci dicono, nei vari casi, 
quale aula deve spostarsi).

\begin{center}
	\psset{unit=.8ex}
	\pspicture(-25,-32)(72,14)
	
	\rput[l](-2,12){\textbf{Prima\ldots}}
	\rput[l](38,12){\textbf{Dopo\ldots}}
	\rput(-20,2.5){$n=3$}
	
	\rput(0,0){
	\rput(-20,2.5){$n=3$}
	\aulagialla{0,0}{1}
	\aula{4,0}{2}
	\aula{0,5}{3}}
	
	\rput(40,0){
	\aulagialla{4,5}{1}
	\aula{4,0}{2}
	\aula{0,5}{3}}
	
	\rput(0,-15){
	\rput(-20,2.5){$n=8$}
	\aulagialla{0,0}{6}
	\aula{4,0}{7}
	\aula{8,0}{5}
	\aula{0,5}{1}
	\aula{4,5}{2}
	\aula{8,5}{3}
	\aula{12,5}{4}
	\aula{16,5}{8}
	}
	
	\rput(40,-15){
	\aulagialla{20,5}{6}
	\aula{4,0}{7}
	\aula{8,0}{5}
	\aula{0,5}{1}
	\aula{4,5}{2}
	\aula{8,5}{3}
	\aula{12,5}{4}
	\aula{16,5}{8}
	}
	
	\rput(0,-30){
	\rput(-20,2.5){$n=11$}
	\aulagialla{0,0}{11}
	\aula{4,0}{4}
	\aula{8,0}{8}
	\aula{12,0}{10}
	\aula{0,5}{1}
	\aula{4,5}{2}
	\aula{8,5}{3}
	\aula{12,5}{5}
	\aula{16,5}{6}
	\aula{20,5}{7}
	\aula{24,5}{9}
	}
	
	\rput(40,-30){
	\aulagialla{28,5}{11}
	\aula{4,0}{4}
	\aula{8,0}{8}
	\aula{12,0}{10}
	\aula{0,5}{1}
	\aula{4,5}{2}
	\aula{8,5}{3}
	\aula{12,5}{5}
	\aula{16,5}{6}
	\aula{20,5}{7}
	\aula{24,5}{9}
	}
	
	\endpspicture
\end{center}
}

\renewcommand{\SoluzioniAlternative}{
}

\renewcommand{\Osservazioni}{
}

\renewcommand{\Commenti}{
}